\chapter{Conclusion}\label{chap:conclusion}

This dissertation has investigated the Multilevel Monte Carlo 
method for two parabolic SPDEs: the Stochastic Heat Equation and the 
Dean-Kawasaki equation. The primary objectives were to quantify the 
computational savings achievable over standard Monte Carlo 
(MC), to evaluate the effect of different coupling strategies, 
and to explore the role of parallelisation in accelerating practical 
implementations. 

Across both equations, our results demonstrate the substantial computational 
benefits of MLMC. In every case, MLMC achieved marked 
reductions in cost compared with MC, in line with 
the complexity bounds of the MLMC theorem. 

For the SHE, the impact of coupling mechanisms was 
particularly striking. The Nearest Neighbour and Central 
Coupling methods exhibited similar, 
suboptimal variance 
decay rates, while the Finite Element coupling method 
achieved a faster variance decay. This, in 
turn, enabled the MLMC algorithm to yield 
further improvements over the MC method and 
highlighted the importance of coupling methods 
in constructing effective MLMC estimators.
We also found that, contrary to expectation, 
the Central Coupling method offered no 
material improvement over the Nearest Neighbour method. The 
two methods often performed identically, particularly 
for smaller RMSEs. 

For the DK equation, performance across coupling 
methods was more uniform. All coupling 
methods achieved a variance decay rate of $\beta \approx 2$, 
resulting in broadly similar performance 
improvements and equivalent cost scalings. 

Parallelisation provided an additional improvement in performance.
Even a naive OpenMP implementation, parallelising only over 
independent samples, produced substantial reductions in wall-clock time. 
These results underscore the natural compatibility 
of MLMC with parallel computing. 

Taken together, these findings address the three central 
questions posed at the outset:

\begin{itemize}
    \item \textbf{Cost savings: } MLMC achieves very large computational 
            gains over MC for both SPDEs studied.
    \item \textbf{Coupling Mechanisms: } Coupling plays a decisive role 
            for the SHE, with the Finite Element method outperforming
            the simpler schemes.
            For the DK equation, the Nearest Neighbour and Central Coupling 
            methods behave equivalently, and the Finite Element method 
            provides no further advantage.
    \item \textbf{Parallelism: } Even simple parallelisation across trials offers a significant 
            additional speedup, most pronounced in the computationally simplest method, but still
            averaging $4\times$ speedups in the lowest average improvement.
\end{itemize}

\bigskip

\textbf{Future Work}

Several natural extensions present themselves. Further theoretical 
analysis to explain why the FE method provides improved correlations 
for the SHE compared to the NN and CC methods, while not yielding any 
improvements over them for the DK equation, is needed. We also 
wish to derive why the FE method achieves the variance decay rates 
observed for the Stochastic Heat Equation cases.
Similarly, exploring alternative coupling methods, such as 
the Fourier based coupling method used by Cornalba and Fischer 
\cite{cornalba2025multilevel}, would be valuable, applying them 
to other classes of SPDEs. Finally, extending this work 
to high-performance implementations that exploit 
the parallelism over levels and within samples could 
further demonstrate the scalability of MLMC in large-scale 
stochastic simulation.

In conclusion, this dissertation has demonstrated 
both the theoretical and practical power of MLMC in reducing the 
cost of computing expectations for SPDEs. We demonstrated the 
importance of coupling mechanisms and introduced a coupling 
method that yielded the optimal cost rate for the Stochastic 
Heat Equation.
We also found that a method which perfectly 
reuses fine level noise is functionally equivalent to a 
method which discards some. Our results highlight that MLMC 
can achieve significant performance improvements when applied 
to Stochastic PDEs, which can be further enhanced by 
taking advantage of its parallel nature.

\textit{The code used to conduct this investigation and to write this report 
can be found at the following: \url{https://github.com/intimantripp/mlmc_for_spdes}}

